\section{Testing and Verification}\label{sec:6}

This section describes how the implementation of Section~\ref{sec:5} was verified: what each check tests, how it tests it and what it found. The design that the checks protect is given in Section~\ref{sec:4}, and the evaluation results are presented in Section~\ref{sec:7}. Verification in TraceAdvisor takes three forms. Automated tests run the shared modules and the Alibaba GPU 2020 scripts \cite{ref19} on small constructed inputs whose correct outputs are known; checks inside the Google ClusterData 2019 notebooks \cite{ref7} and the Alibaba scripts stop a stage when a condition fails on the real data; and the saved outputs of every stage record counts and checksums that later stages verify. The full test suite was run on 30 October 2026 at the final commit in the locked environment: 699 tests passed and one was skipped. Of these, 658 cover the two datasets in scope; the remainder belong to adapters for other traces that are not evaluated in this report. The skipped test loads the full official Alibaba tables and runs only when they are downloaded and explicitly enabled. Table~\ref{tab:6.1} summarizes the checks by stage.

\begin{table*}[!t]
\centering
\footnotesize
\caption{Verification by pipeline stage}\label{tab:6.1}
\begin{tabular}{>{\raggedright\arraybackslash}p{0.385\linewidth}>{\raggedright\arraybackslash}p{0.067\linewidth}>{\raggedright\arraybackslash}p{0.457\linewidth}}
\hline
\textbf{Stage} & \textbf{Automated tests} & \textbf{Checks on the real data} \\ \hline
Canonical schema and adapter framework & 147 & Alibaba accepted and rejected records validated against the schema at build time \\
Google canonical contract & 5 & Saved Google occurrence tables mapped to the schema \\
Alibaba source loading, adapter and canonical build & 259 & Source checksums; zero difference between source rows and accepted plus rejected records; 461,900 eligible role occurrences \\
Splits, histories and targets & 86 & Timestamp groups kept within one split; Google histories and profiles recomputed independently for 4,879 occurrences \\
Baselines & 28 & Training and validation histories unchanged when test rows are added \\
ML preparation, selection and final test & 48 & Test targets never read before the decision manifest; one recorded final-test run per dataset \\
Hybrid rule & 10 & No row with a missing safeguard on Alibaba \\
Metrics and evaluator & 75 & Google Notebook 06 reproduces the Notebook 04 baseline metrics exactly \\
\hline
\end{tabular}
\end{table*}

\subsection{Data and Pipeline Correctness Checks}\label{sec:6.1}

The canonical schema is the contract between the dataset-specific stages and the shared ones, so it is tested most heavily. The schema tests construct records that violate one rule at a time, such as a missing identifier, a timestamp in the wrong unit, a resource value with an undeclared meaning or a validity label that contradicts the value, and confirm that each is rejected with a specific message while valid records pass unchanged. The adapter framework tests add the rules that span records: unique keys, consistent recurrence keys, a rejected record for every excluded input and a valid parent for every child record.

For Google, the notebooks verify each reconstruction step on the data itself. Notebook 00 found 405,894 observations with no duplicate observation keys, and Notebook 01 found no observation whose peak usage was below its average usage when both were recorded. Notebook 02 reconstructed 253,585 executions with no negative runtime and no missing execution or recurrence key. Because an occurrence combines several executions, Notebook 03 recomputes each occurrence's peaks, runtime, requests and history statistics by a separate, simpler calculation and compares them with the pipeline's values; all 4,879 occurrences and 87,822 history statistics agreed, and statistics that are undefined for a workload's first runs remained missing rather than zero. These checks establish that the reconstruction is internally consistent; they cannot establish that every reconstructed recurrence corresponds to the workload a user submitted, since the trace records no such ground truth. The Google canonical contract tests map the saved Google occurrence tables to the shared schema, which confirms that the notebook outputs satisfy the same contract as the Alibaba adapter even though Google has no production adapter.

For Alibaba, verification starts at the source files. The loader checks each official table against its recorded checksum, size, line count and header before any row is parsed, and the adapter tests cover the joins between jobs, task roles and worker instances, the exclusion reasons and the aggregation rules. Two aggregation rules are tested directly because an error in either would change every result without failing elsewhere: a role's resource value is the maximum over its covered workers, never their sum, and a missing value on a covered worker makes the role's value missing rather than zero. The canonical build reconciles its counts on the real data, and the saved report shows no difference between the number of source rows and the number of accepted plus rejected records. Of the 461,900 eligible role occurrences, 1,224 have no CPU measurement but a valid memory measurement, which is why CPU and memory validity are recorded separately.

Every Alibaba stage writes a \texttt{SHA256SUMS} file with its outputs, and each later stage verifies the files it reads against that list before using them; the target builder also checks the row counts and column order of the canonical tables. A file altered, truncated or regenerated after a stage completed therefore stops the next stage. A planned independent verification, a separate script that re-validates every artifact and reproduces the pipeline in a fresh environment, was not completed; the evidence for the Alibaba pipeline is therefore the per-stage verification above, the end-to-end tests of Section~\ref{sec:6.2} and the logged final run.

\subsection{Temporal Leakage and Split Integrity}\label{sec:6.2}

A recommendation for a run may use only information available when that run starts, so the checks in this section test that no future information reaches a feature, a baseline, a model or a selection decision. Both pipelines split their rows chronologically into approximately 70\%, 15\% and 15\% for training, validation and test. The split tests construct rows that share a timestamp and confirm that they are never separated, that the partitions contain distinct identifiers and that every training row starts before every validation row, and every validation row before every test row. On Alibaba the target builder additionally refuses an occurrence whose roles would fall into different splits. A chronological split orders rows by start time; it does not prevent a training run from still executing when a validation run starts, which is why history eligibility is checked separately.

History eligibility is tested against independent calculations. A run may enter another run's history only if it completed at or before that run started. Notebook 03 checks this rule on every Google occurrence, and the Alibaba target builder is compared with a brute-force calculation that, for each row, scans all earlier rows and applies the rule directly. Further tests cover the boundary cases: a predecessor that completes exactly at the current start is eligible, an overlapping unfinished predecessor is not, a zero-duration predecessor with the same submission time is ordered deterministically, and CPU and memory histories are built independently so that a missing CPU value does not remove a memory value. A separate test changes every measurement in a later split and confirms that no feature or target of an earlier split changes. The Alibaba baseline script repeats this check on the real data: it computes the training and validation history statistics with and without the test rows and stops if they differ.

Target isolation differs between the datasets. Alibaba assigns a target only when the next eligible occurrence is in the same split; otherwise the target is left missing with the status \texttt{next\_occurrence\_in\_other\_split}, and tests confirm that histories may draw on earlier splits while targets may not. Google builds its next-occurrence targets before splitting, so the check measures the effect rather than preventing it: 104 of the 329 validation occurrences have a target that starts in the test period. These validation targets are therefore not independent of the test period, and they were used to select the Google baseline and models. Notebook 06 rescored validation on the 225 occurrences whose targets remain in the validation period. The selected baseline and the order of the core policies were unchanged (Section~\ref{sec:7}), but this sensitivity analysis measures the limitation rather than removing it.

Model inputs are checked against a list of approved features. Google audits its feature manifest against the known outcome columns and trains with folds that are purged at the start time of each training label's target, and Alibaba validates the training, validation and test feature tables against its list of leaky columns. Google's purge uses the start time of the target occurrence, which does not establish that the target's complete usage was available by the cutoff.

The final test is protected by checks that run before the test targets are read. On Google, Notebook 06 freezes every decision into a decision manifest in one run and evaluates the test set in a second run. The evaluation run refuses to start unless the manifest is final, the current commit equals the frozen commit, the only uncommitted change is the notebook itself, the notebook's code hash and every hashed input still match, and no final-test result already exists. The recorded run evaluated the 259 test occurrences with both targets once. On Alibaba, the baseline script writes its selection record before it opens the test table, and then reads only an allowed list of test columns that excludes the targets. The ML script never opens the test table: its end-to-end test records every table the script reads and fails if the test table is among them. The final-test script then runs its checks in a fixed order. It verifies the checksums of the target, baseline and ML outputs against the manifest, including the hash of the test table, and checks that the manifest's fields are complete. It confirms that the manifest was built from a clean working tree at the current commit. It then reads its run log, which is keyed by the hash of the decision manifest and kept with the locked models rather than in the output folder, and refuses a second run for the same manifest whatever output folder is requested. Only then does it append the run to the log and open the test table, so a run that fails after reading the test set is still counted. A rerun is allowed only with a documented defect reason; it is logged, and the previous output is preserved. The ML script in turn refuses to replace a lock that has already been evaluated on the test set. Each of these refusals has its own test, which also confirms that the test table was not opened, and an end-to-end test runs the baseline, ML and final-test scripts in sequence, then confirms that a second final-test run and a replacement of the lock are both refused. On the real data the run log contains a single entry, with no prior run and no defect reason.

\subsection{Model and Baseline Validation}\label{sec:6.3}

The baseline tests verify the history statistics, the fallback rule and the selection. They confirm that each row's history count matches the number of eligible predecessors and that a mismatch is rejected. They also confirm that a row with fewer than \texttt{min\_runs} usable runs for a resource falls back to its request for that resource only, that every candidate is scored on the same population and that the ranking follows the order of Section~\ref{sec:4.2.1}, including that Request is never selected on Alibaba. On Google, Notebook 04 checks that Request gives identical results at every \texttt{min\_runs}, since it does not use history, and Notebook 06 recomputes the baseline validation metrics and confirms that they equal those of Notebook 04 exactly.

The model checks verify that selection uses validation data only and that the selected models are the ones evaluated on test. Google tunes each candidate with purged time-ordered folds on the training set and selects per resource by validation RMSE on the logarithmic scale; the notebook records the full tuning results and the selection in its metadata, and both selected models are random forests. Alibaba compares fixed configurations of a decision tree, a random forest, histogram gradient boosting and XGBoost, together with a constant mean predictor, and selects per resource by validation RMSE in the original units. Its tests confirm that all candidates fit and predict, that the constant predictor is reported but never selected, and that selection fails if no other candidate remains. They also confirm that the models refitted on training and validation, which are locked for the test, differ from the models fitted on training alone, which are used for the validation comparison. Histogram gradient boosting was selected for both resources. The final-test script loads the locked model files, checks their hashes against the manifest and does not refit them.

The Hybrid rule takes the larger of the ML prediction and the safeguard, so its tests check the property on which its interpretation depends: the recommendation is never below either input. They also check that a missing safeguard falls back to the ML prediction, and that negative or non-finite inputs and misaligned rows are rejected. Google uses P95 as its safeguard, and Alibaba the selected baseline, P99 with \texttt{min\_runs} of 10. The count of rows on which the safeguard was missing is recorded for each run: on Google one CPU and one memory row fell back to the ML prediction on validation, and on Alibaba no row did, on validation or on test. The saved final results are consistent with the property: Hybrid has fewer violations than its safeguard on every resource on both datasets (Section~\ref{sec:7}).

\subsection{Replay Evaluation Metrics}\label{sec:6.4}

Both pipelines score their policies with the same evaluator, which compares each recommendation with the recorded usage of the next occurrence and is defined in Section~\ref{sec:4.3}. The evaluator tests use small hand-computed cases whose correct metrics are known, and contract tests check its behavior under each dataset's capability profile. Replay measures whether a recommendation would have covered the recorded usage; it does not simulate scheduling or measure application performance.

\subsubsection{Violation and Utilization}\label{sec:6.4.1}

The violation tests fix the boundary cases of the definition. Usage equal to the recommendation is not a violation, usage strictly above it is, and a recommendation of zero is a valid comparison that is violated by any positive usage. If either value is missing or non-finite the result is unknown rather than a non-violation. The combined Google metric has three states: violated if either resource is violated, not violated only if both are known not to be, and unknown otherwise. The utilization tests confirm that utilization is the ratio of usage to recommendation capped at one, and that it is undefined for a recommendation of zero or near zero rather than infinite. The contract tests confirm that the Alibaba profile labels CPU results as lifetime-average usage, identifies memory as the primary resource and suppresses the combined metric, because a combined rate would mix an average with a peak; the saved Alibaba results carry these labels.

\subsubsection{Resource Waste}\label{sec:6.4.2}

The waste tests confirm that waste is the amount by which the recommendation exceeds usage and zero otherwise, and that shortfall is the amount by which usage exceeds the recommendation, so that exactly one of the two is positive on any violated or over-provisioned row. Both are unknown when either input is missing, and the tests confirm that unknown values are excluded from the averages rather than counted as zero. Waste is averaged per row without weighting by run duration, and is reported separately for CPU and memory in each dataset's own units, so it is not compared across the datasets.

\subsubsection{Valid-Row Denominators and Missing-Value Handling}\label{sec:6.4.3}

Each metric is averaged over the rows on which it is defined, so different metrics can have different denominators within the same policy. The denominator tests confirm that the evaluator returns these counts and that a row with an unknown outcome is not counted as a non-violation. On Google, the population is the occurrences with both targets, but a policy can still lack a recommendation on some of them: of the 329 validation occurrences, Request has a value on 326, and of the 259 test occurrences, 257 have a Request CPU and memory violation outcome and 254 a CPU utilization value. On Alibaba, a row enters the population if either target is available, and each resource is scored on its own rows, so that a missing CPU target does not remove a valid memory target: validation has 42,069 memory and 41,913 CPU outcomes, and test 44,053 and 44,045. The Alibaba ML features fill missing history values with zero and add an indicator for each, but this applies only to model inputs; a missing target is never scored as zero usage. The saved comparison tables export the violation and utilization counts but not the waste and shortfall counts, which the profiled evaluator computes but the scripts do not write.
